\section{AI Adoption: del Mr. T Demo al Category-Defining Workflow}

Al final de la entrevista, McKinnon compara la IA con las primeras aplicaciones del iPhone: la capacidad del hardware o del modelo importa, sin duda, pero al principio muchas aplicaciones solo envolvían la nueva capacidad como novelty; lo que de verdad transformó la industria fue el workflow al estilo de Uber, que llegó después. Este juicio no afirma que la IA carezca de valor, sino que advierte a las empresas incumbent y a los equipos emprendedores que distingan entre capability shock, demo, repeatable workflow y organization redesign.

\subsection{Capacidad técnica e inercia organizacional}

Las grandes empresas tienen clientes, datos y procesos, pero también role, budget, support model e inercia cultural. Aunque un AI agent pueda completar parte de una tarea de support, operations o security, la organización todavía necesita reescribir accountability, quality review, exception handling, identity delegation y performance metric. Si solo se coloca un cuadro de chat delante del proceso anterior, el sistema puede terminar con una «Mr. T app» sin haber cambiado la outcome economics.

\lecturefigure{16-ai-adoption-curve.png}{Por lo general, la IA pasa primero por un capability shock y por novelty apps; solo después aparecen el category-defining workflow y el rediseño organizacional.}{Subtítulos locales 36:00--39:02; redibujo conceptual.}

\begin{knowledgebox}{Lectura de la figura: cuatro preguntas para identificar un killer workflow}
¿Completa una tarea que antes no podía realizarse de forma económica, en lugar de solo cambiar la UI? ¿Su uso genera compounding data/feedback? ¿Las fallas tienen un owner y una recovery claros? ¿La organización está dispuesta a cambiar role, process, control y pricing? Si todas las respuestas son negativas, la aplicación probablemente siga siendo una demo interesante; solo si cambian a la vez el workflow, la distribution y la governance puede la tecnología formar una nueva category.
\end{knowledgebox}

\teachervoice{McKinnon no adopta un relato de AI doom ni afirma que el killer use case ya esté definido. Compara la etapa actual con la aplicación de Mr. T de los primeros tiempos del iPhone, que decía frases ingeniosas: todos saben que la plataforma importa, pero «un workflow completamente nuevo que nace de una nueva capacidad, como Uber» todavía está en formación.}

\begin{warningbox}{La incumbent inertia es a la vez un riesgo y un margen de seguridad}
Los procesos anteriores frenan la automatización valiosa, pero también pueden impedir por un tiempo que un agent sin gobernanza obtenga permisos de producción. El rediseño organizacional debe, al mismo tiempo, eliminar los handoff ineficaces y añadir delegation, review, audit y fallback; «AI-first» no puede servir de pretexto para eludir identity, change management e incident response.
\end{warningbox}

\subsection{Resumen de la sección}

El valor de la IA va de la capability al workflow y, después, a la organization. En este trayecto, la identity infrastructure no es un componente de inicio de sesión que se añade al final, sino la condición previa para la agent accountability, el least privilege y el cross-tool audit.
